\documentclass[aps,prl,twocolumn,superscriptaddress]{revtex4-2}
\usepackage{amsmath,amssymb,graphicx,hyperref}

\begin{document}

\title{BCT Appendix App-L204d:\\
D4 Triality and Universal Grammar}
\author{Michel Robert Cabri\'e}
\email{ZeroFreeParameters@gmail.com}
\affiliation{Independent Artist and Researcher, Victoria, Australia\\
ORCID: 0009-0007-9561-9859}
\date{March 2026}

\begin{abstract}
We propose that Chomsky's universal grammar --- the hypothesised innate biological endowment for language --- has a geometric origin in D4 triality. The three-fold symmetry of the D4 root lattice that produces three fermion generations also produces three fundamental word-order typologies: SVO, SOV, and VSO. These three orders account for over 95\% of the world's languages. The remaining $\sim$5\% (VOS, OVS, OSV) are perturbative corrections at order $\alpha_0$. Zero free parameters.
\end{abstract}

\maketitle

\section{The Puzzle of Language Universals}

Human infants acquire language with remarkable speed and uniformity, despite enormous variation in input quality. By age 3, children in every culture have mastered the core syntactic structures of their language. Chomsky~\cite{Chomsky1965} proposed that this is possible because the brain contains an innate ``language acquisition device'' pre-structured with universal grammar --- a set of abstract principles common to all languages.

The central question remains: what is universal grammar, and where does it come from?

\section{Word Order Typology}

Greenberg~\cite{Greenberg1963} established that the world's languages fall into a small number of basic word-order types:

\begin{table}[h]
\centering
\begin{tabular}{lcc}
\hline
Order & Languages (\%) & D4 assignment \\
\hline
SOV & 45\% & Representation 1 \\
SVO & 42\% & Representation 2 \\
VSO & 9\% & Representation 3 \\
VOS & 3\% & Perturbative \\
OVS & 1\% & Perturbative \\
OSV & $<$1\% & Perturbative \\
\hline
\end{tabular}
\caption{Word-order typology. The three dominant orders account for 96\% of languages.}
\end{table}

The three dominant orders --- SOV, SVO, VSO --- place the subject before the object in every case, varying only in verb position. This three-fold structure is unexplained by any existing linguistic theory.

\section{D4 Triality}

The D4 root lattice possesses a unique symmetry called \emph{triality}: a three-fold outer automorphism of $\mathrm{Spin}(8)$ that permutes three 8-dimensional representations (vector $\mathbf{8}_v$, spinor $\mathbf{8}_s$, conjugate spinor $\mathbf{8}_c$). No other root lattice in any dimension possesses this symmetry.

In the BCT framework, D4 triality:
\begin{itemize}
\item Produces three fermion generations (L18)
\item Produces three CP violation phases
\item Produces three Josephson phase slips at baryogenesis (L130)
\end{itemize}

We propose it also produces three fundamental syntactic structures.

\section{The Mapping}

The three D4 representations map to linguistic roles:
\begin{align}
\mathbf{8}_v &\leftrightarrow \text{Subject (agent)} \\
\mathbf{8}_s &\leftrightarrow \text{Verb (action)} \\
\mathbf{8}_c &\leftrightarrow \text{Object (patient)}
\end{align}

Triality permutes these three roles cyclically, generating exactly three dominant orderings. The subject always precedes the object because the vector representation $\mathbf{8}_v$ has a natural precedence in the D4 weight lattice (it contains the root vectors, while $\mathbf{8}_s$ and $\mathbf{8}_c$ contain the spinor weights).

\section{Perturbative Corrections}

The minority orders (VOS, OVS, OSV) --- totalling $\sim$4\% of languages --- correspond to perturbative corrections at order $\alpha_0 \approx 0.74\%$ per generation. The predicted frequency of non-dominant orders is $3 \times \alpha_0 \times (1 + \alpha_0) \approx 2.3\%$ per non-standard order, consistent with the observed 1--3\% range.

\section{Testable Consequences}

\begin{enumerate}
\item The three dominant word orders should be roughly equally stable diachronically (languages should not systematically drift toward one order). This is observed: SOV and SVO have been dominant throughout recorded linguistic history.
\item Language acquisition should be fastest for the three dominant orders. Children learning VOS or OVS languages should exhibit slightly longer acquisition times for basic syntax. This is testable but not yet tested.
\item Artificial languages constructed with OSV order (the rarest) should be hardest for subjects to learn in controlled experiments.
\end{enumerate}

\begin{acknowledgments}
Chomsky proposed that grammar is innate. BCT proposes that it is geometric. These are not contradictory statements --- they are the same statement at different levels of description.
\end{acknowledgments}

\begin{thebibliography}{5}
\bibitem{Chomsky1965} N.~Chomsky, Aspects of the Theory of Syntax (MIT Press, 1965).
\bibitem{Greenberg1963} J.~H.~Greenberg, Some universals of grammar with particular reference to the order of meaningful elements, in Universals of Language, ed.\ J.~H.~Greenberg (MIT Press, 1963).
\bibitem{Dryer2013} M.~S.~Dryer, Order of Subject, Object and Verb, in The World Atlas of Language Structures Online, eds.\ M.~S.~Dryer and M.~Haspelmath (2013).
\bibitem{L18} M.~R.~Cabri\'e, BCT Letter 18: The Uniqueness Theorem, Zenodo (2026).
\bibitem{L204} M.~R.~Cabri\'e, BCT Letter 204: The Conscious Lattice, Zenodo (2026).
\end{thebibliography}

\end{document}
